\begin{ouverture}[derivation]{5}{Convexité et inégalités}{Analyse et répartition}{Sécantes · Tangentes · Inflexion · Lorenz}
\objectif{Reconnaître convexité et concavité sur un graphique.}
\objectif{Utiliser le signe de la dérivée seconde.}
\objectif{Justifier une inégalité par une tangente.}
\objectif{Lire une courbe de Lorenz et comparer des répartitions.}
\prerequis{Dérivation, variations et pourcentages cumulés.}
\end{ouverture}
\essentiel

\essSection{1}{La position des cordes}
\begin{definition}
Une fonction est \textbf{convexe} sur un intervalle si sa courbe est située sous chaque segment reliant deux de ses points. Elle est \textbf{concave} si sa courbe est au-dessus de ces segments.
\end{definition}
Une fonction peut être croissante et concave : les valeurs augmentent alors que les pentes diminuent. La convexité décrit la variation des pentes, pas le signe de la fonction ni son sens de variation.
\figurelivre{c05-convexite.pdf}{La parabole est sous la corde et au-dessus de ses tangentes. Pour le logarithme, les positions sont inversées : il est concave sur $]0,+\infty[$.}
\essSection{2}{Deux critères admis}
\begin{propriete}
Pour $f$ dérivable, $f$ est convexe si et seulement si $f\,\!^{\prime}$ est croissante ; sa courbe est alors au-dessus de toutes ses tangentes.
Si $f$ est deux fois dérivable, $f\,\!^{\prime\prime}\geq0$ assure la convexité et $f\,\!^{\prime\prime}\leq0$ la concavité.
\end{propriete}
\begin{definition}[Inflexion]
Un point d'inflexion est un point de la courbe où la convexité change. Dans les cas étudiés ici, un changement de signe de $f\,\!^{\prime\prime}$ permet de le reconnaître.
Une dérivée seconde nulle ne suffit pas : $x^4$ reste convexe de part et d'autre de $0$.
\end{definition}
\begin{exemple}[Une courbe en S]
Pour $f(x)=x^3$, $f\,\!^{\prime\prime}(x)=6x$ : concavité sur $]-\infty,0]$, convexité sur $[0,+\infty[$ et inflexion en $(0,0)$.
\end{exemple}
\essSection{3}{Obtenir une inégalité}
L'exponentielle est convexe et sa tangente en $0$ est $y=1+x$. On obtient, pour tout réel $x$, $\mathrm e^x\geq1+x$.
Le logarithme est concave et sa tangente en $1$ est $y=x-1$. Pour $x>0$, $\ln x\leq x-1$.
\begin{savoirfaire}[Étudier une convexité]
Calculer $f\,\!^{\prime}$ puis $f\,\!^{\prime\prime}$ ; factoriser et étudier le signe sur le domaine ; annoncer les intervalles ; vérifier un changement effectif de signe avant de conclure à une inflexion ; calculer les coordonnées du point.
\end{savoirfaire}
\essSection{4}{Classer une population pour lire une répartition}
On classe les personnes par richesse croissante. Une courbe de Lorenz associe à une proportion cumulée $x$ de population la proportion cumulée $L(x)$ de richesse possédée.
Dans les modèles usuels du thème, $L(0)=0$, $L(1)=1$, $L$ est croissante et convexe sur $[0,1]$, et $L(x)\leq x$.
La diagonale $L(x)=x$ correspond à l'égalité parfaite.
\figurelivre{c05-lorenz.pdf}{Les courbes $L(x)=x^2$ et $L(x)=x^3$ sont des modèles fictifs. À $x=0{,}5$, la première attribue $25\,\%$ de la richesse à la moitié la moins riche ; la seconde, $12{,}5\,\%$.}
\begin{exemple}[Ne pas inverser les axes]
Si $L(0{,}8)=0{,}6$, les $80\,\%$ les moins riches possèdent $60\,\%$ de la richesse. Les $20\,\%$ les plus riches en possèdent donc $40\,\%$. Cela ne signifie pas que $60\,\%$ de la population possède $80\,\%$ de la richesse.
\end{exemple}
\essSection{5}{Quantiles et dispersion}
La médiane partage la série ordonnée en deux groupes. $Q_1$ et $Q_3$ délimitent les quartiles ; $D_1$ et $D_9$ les déciles.
Le rapport interdécile $D_9/D_1$, lorsque $D_1>0$, compare deux seuils ; il ne donne pas le rapport des richesses totales de deux groupes.
Pour une petite série discrète, on annonce la convention de quantile employée. Au chapitre 7, le calcul intégral donnera un autre indicateur, l'indice de Gini.

\exercices
\markright{Prendre confiance}
\begingroup\renewcommand{\otBesoinPlace}[1]{}\begin{minipage}{\linewidth}
\rubrique{Prendre confiance}
\begin{exercice}[1]{Signe de la dérivée seconde}

\begin{questions}
\item Étudier la convexité de $x^2+3x$.
\item Étudier celle de $-x^2$.
\item Étudier celle de $\mathrm e^x$.
\end{questions}
\end{exercice}
\begin{exercice}[1]{Logarithme et racine}

\begin{questions}
\item Calculer la dérivée seconde de $\ln x$.
\item Quelle est sa convexité ?
\item Calculer la dérivée seconde de $\sqrt x$ pour $x>0$.
\end{questions}
\end{exercice}
\begin{exercice}[1]{Inflexion}
$f(x)=x^3-3x$.
\begin{questions}
\item Calculer $f\,\!^{\prime\prime}$.
\item Donner les intervalles de convexité et concavité.
\item Donner le point d’inflexion.
\end{questions}
\end{exercice}
\begin{exercice}[1]{Lire Lorenz}
Une courbe fictive vérifie $L(0{,}4)=0{,}1$ et $L(0{,}9)=0{,}7$.
\begin{questions}
\item Interpréter la première valeur.
\item Quelle richesse possèdent les $10\,\%$ les plus riches ?
\item Quelle richesse possède le groupe compris entre les proportions $0{,}4$ et $0{,}9$ ?
\end{questions}
\end{exercice}
\end{minipage}\endgroup
\clearpage
\markright{Appliquer et relier}
\begingroup\renewcommand{\otBesoinPlace}[1]{}\begin{minipage}{\linewidth}
\rubrique{Appliquer et relier}
\begin{exercice}[2]{Une exponentielle amortie}
$f(x)=x\mathrm e^{-x}$ sur $[0,+\infty[$.
\begin{questions}
\item Calculer $f\,\!^{\prime\prime}$.
\item Étudier la convexité.
\item Donner le point d’inflexion.
\end{questions}
\end{exercice}
\begin{exercice}[2]{Une tangente utile}

\begin{questions}
\item Donner la tangente à $\mathrm e^x$ en zéro.
\item En déduire une minoration de $\mathrm e^{0{,}2}$.
\item Donner une majoration de $\ln(1{,}2)$ par la tangente du logarithme en $1$.
\end{questions}
\end{exercice}
\begin{exercice}[2]{Comparer deux répartitions}
$L_A(x)=x^2$ et $L_B(x)=x^3$ sur $[0,1]$.
\begin{questions}
\item Vérifier les extrémités et la croissance.
\item Comparer les deux courbes.
\item Quel modèle est le plus éloigné de l’égalité au sens de Lorenz ?
\end{questions}
\end{exercice}
\begin{exercice}[2]{Déciles}
Dans une série fictive de revenus, $D_1=900$, la médiane vaut $1\,800$ et $D_9=3\,600$ euros.
\begin{questions}
\item Calculer le rapport interdécile.
\item Peut-on affirmer que les $10\,\%$ les plus riches possèdent quatre fois la richesse des $10\,\%$ les moins riches ?
\item Quel renseignement manque pour construire une courbe de Lorenz ?
\end{questions}
\end{exercice}
\end{minipage}\endgroup
\clearpage
\markright{Chercher et argumenter}
\begingroup\renewcommand{\otBesoinPlace}[1]{}\begin{minipage}{\linewidth}
\rubrique{Chercher et argumenter}
\begin{exercice}[3]{Zéro sans inflexion}
$f(x)=x^4$.
\begin{questions}
\item Calculer la dérivée seconde.
\item Que se passe-t-il en zéro ?
\item Formuler une règle correcte pour repérer un point d’inflexion dans les exemples deux fois dérivables.
\end{questions}
\end{exercice}
\begin{exercice}[3]{Un modèle admissible ?}
On propose $L(x)=2x-x^2$ pour une courbe de Lorenz.
\begin{questions}
\item Vérifier les extrémités et la croissance.
\item Comparer $L(x)$ à $x$.
\item Expliquer pourquoi ce modèle ne convient pas avec un classement par richesse croissante.
\end{questions}
\end{exercice}
\begin{exercice}[3]{Convexité sans croissance}
$f(x)=\mathrm e^{-x}$ sur $\R$.
\begin{questions}
\item Calculer les deux premières dérivées.
\item Donner le sens de variation et la convexité.
\item Réfuter « toute fonction convexe est croissante ».
\end{questions}
\end{exercice}
\begin{exercice}[3]{Une courbe logistique}
$f(x)=1/(1+\mathrm e^{-x})$. On donne $f\,\!^{\prime\prime}(x)=\mathrm e^{-x}(\mathrm e^{-x}-1)/(1+\mathrm e^{-x})^3$.
\begin{questions}
\item Étudier le signe de la dérivée seconde.
\item Déterminer l’inflexion.
\item Interpréter le changement de convexité pour une croissance avec saturation.
\end{questions}
\end{exercice}
\end{minipage}\endgroup
